\scanpage{12}
\deleted{fixe, mettons $\mathbb C$ ici) sont déterminés à ext. alg. près. Nous allons examiner \unsure{la|le}

\uline{Conjecture} \inserted{\uline{de transcendance}} Si $k$ est \unsure{quasi-monnaire|quasi-monétaire|quasi-unitaire} pour $M$, alors $\operatorname{deg\,tr}(k(g_0))=\dim G^{\unsure{u|n|\prime\prime}}\;(=\dim G-1)$, $\deleted{\unsure{\xi}}g_0$ \inserted{(localité de $g_0$)} est un point \uline{maximal} de $G^u$ (donc gén. si \unsure{$G$ de} $G^u$ connexe).
\margin{\unsure{caracté|caract.} les courbes elliptiques à invariant transcendant et $\tau$ algébrique\ldots\unsure{?!}}% slanted

\uline{Corollaire} \inserted{Si $G$ connexe} $G^u$ est le $k$-groupe algébrique engendré par $g_0$. Donc $G$ est le $k$-groupe algébrique engendré par $g_0$ et $j(G_m)$ (les homothéties). Enfin $G$ est engendré sur $k$ par $j$ si $G^u$ est $\tfrac12$ simple et $\pi$ non algébrique sur $k$ \deleted{\unsure{connexe}}\ldots
\margin{et cas général ? \deleted{\unsure{gr.}\;$\pi$ \unsure{non alg. sur $k$ ?}}}

Contient le th. de \deleted{\unsure{D}}Schneider, pour les courbes elliptiques sans mult. complexe. En fait, il signifie que le degré de transc. du système $(\underset{\substack{\parallel\\\omega'_1}}{\zeta^{-1}\omega_1},\underset{\substack{\parallel\\\omega'_2}}{\zeta^{-1}\omega_2},\underset{\substack{\parallel\\\unsure{\eta|y}'_1}}{\zeta^{-1}\eta_1},\underset{\substack{\parallel\\\eta'_2}}{\zeta^{\unsure{-1|1}}\eta_2})$ sur $Q$ est 3, \deleted{\unsure{si}} (la seule relation \deleted{\unsure{étant|est}} $\omega'_1\eta'_2-\omega'_2\eta'_1=1$, i.e. l'équation de $G^u=\mathrm{Sl}(2)\subset G=\mathrm{Gl}(\unsure{2|7})$) pour une courbe elliptique sans mult. complexe définie sur un corps de nbs alg. (disons sur $\overline Q$) Si l'invariant $j$ est transcendant, on trouve que $(j,\omega'_1,\omega'_2,\eta'_1,\eta'_2)$ ont le degré de \uline{trans. 4} sur $Q$.}
% archive sheet number: 11
